\section{Hard-edge removal and bulk persistence}
\label{sec:scale-bulk}

We begin with a coarse form of scale localization.  It removes both factor
edges but leaves the ratio $n/m$ otherwise unrestricted.  Define
\begin{align}
 \cP_h(X;W)
 &=\sum_{m,n\geq1}\Lambda(n)\Lambda(mn+h)W(mn/X),
 \label{eq:prime-source-complete}\\
 \cA_h(X;W)
 &=\sum_{m,n\geq1}(\Lambda(n)-1)\Lambda(mn+h)W(mn/X),
 \label{eq:residual-complete}
\end{align}
and let a superscript $\mathrm{bulk}(Y)$ mean that both $m>Y$ and $n>Y$
are imposed.

\subsection{The two complete collapses}

The prime-source form has an even simpler zero-mode collapse than the
residual form.

\begin{proposition}[Complete prime-source collapse]
\label{prop:prime-source-collapse}
For every $A>0$,
\begin{equation}\label{eq:prime-source-complete-asymptotic}
 \cP_h(X;W)
 =X\{(\log X)I_0(W)+I_1(W)\}
 +O_{A,h,W}\!\left(X(\log X)^{-A}\right).
\end{equation}
\end{proposition}

\begin{proof}
Put $r=mn$ and use $\sum_{n\mid r}\Lambda(n)=\log r$.  Then
\[
 \cP_h(X;W)=\sum_{r\geq1}(\log r)\Lambda(r+h)W(r/X).
\]
The fixed-shift prime number theorem and smooth partial summation give
\[
 \sum_r(\log r)\Lambda(r+h)W(r/X)
 =\int_0^\infty (\log u)W(u/X)\,du
 +O_{A,h,W}(X(\log X)^{-A}),
\]
and the integral is the main term in
\eqref{eq:prime-source-complete-asymptotic}.
\end{proof}

The residual form is \eqref{eq:intro-complete-zero-mode}.  Only its leading
coefficient is needed below:
\begin{equation}\label{eq:residual-leading-only}
 \cA_h(X;W)
 =(1-C_h)I_0(W)X\log X+O_{h,W}(X).
\end{equation}

\subsection{An edge ledger}

We record the upper bounds used to remove the two factor edges.  They are
upper-sieve statements; no cancellation between rows is asserted.

\begin{lemma}[One-prime edge]
\label{lem:one-prime-edge}
Uniformly for $2\leq Y\leq X^{1/3}$,
\begin{align}
 &\sum_{q\leq Y}(\Lambda(q)+1)
   \sum_{k\geq1}\Lambda(qk+h)|W(qk/X)|
 \ll_{h,W}X\log(2Y),
 \label{eq:one-prime-q-edge}\\
 &\sum_{q\leq Y}
   \sum_{k\geq1}\Lambda(qk+h)|W(qk/X)|
 \ll_{h,W}X\log(2Y).
 \label{eq:one-prime-k-edge}
\end{align}
\end{lemma}

\begin{proof}
For $(q,h)=1$, the weighted Brun--Titchmarsh inequality and partial
summation give
\begin{equation}\label{eq:weighted-BT-row}
 \sum_{k\geq1}\Lambda(qk+h)|W(qk/X)|
 \ll_{h,W}\frac{X}{\varphi(q)}
\end{equation}
uniformly for $q\leq X^{1/3}$.  The proper prime-power targets contribute
at most $O_W(X^{1/2}(\log X)^2)$ in one row, which is
$O(X/\varphi(q))$ in this range for large $X$: there are
$O_W(X^{1/2})$ powers $p^a\asymp_W X$ with $a\geq2$, and each von
Mangoldt weight is $O(\log X)$.  The elementary estimates
\begin{equation}\label{eq:phi-harmonic-estimates}
 \sum_{q\leq Y}\frac1{\varphi(q)}\ll\log(2Y),
 \qquad
 \sum_{q\leq Y}\frac{\Lambda(q)}{\varphi(q)}\ll\log(2Y)
\end{equation}
then prove the coprime part of both claims.

If $(q,h)>1$ and $qk+h$ is in the support of $\Lambda$, some prime
$p\mid h$ divides $qk+h$.  Hence $qk+h=p^a$.  For each of the finitely
many $p\mid h$, only $O_{h,W}(1)$ powers lie in the fixed-ratio interval
forced by $W(qk/X)$.  Summing the divisor weights over the corresponding
$q\mid p^a-h$ costs $X^\eps$ for every $\eps>0$, which is absorbed by the
right-hand sides.
\end{proof}

The other edge contains two prime weights.  Its standard Selberg-sieve
bound is included with the local-factor calculation because that
dependence on the varying coefficient $m$ is important.

\begin{lemma}[Two-linear-forms edge]
\label{lem:two-linear-forms-edge}
Uniformly for $2\leq Y\leq X^{1/3}$,
\begin{equation}\label{eq:two-linear-forms-edge}
 \sum_{m\leq Y}\sum_{n\geq1}
 \Lambda(n)\Lambda(mn+h)|W(mn/X)|
 \ll_{h,W}X\log(2Y).
\end{equation}
\end{lemma}

\begin{proof}
First suppose $(m,h)=1$ and restrict both von Mangoldt factors to primes.
On an interval $n\asymp_W X/m$, apply the Selberg upper-bound sieve to the
two forms
\[
 n,\qquad mn+h.
\]
If $\nu_p$ is the number of roots of $n(mn+h)$ modulo $p$, then
\[
 \nu_p=
 \begin{cases}
  1,&p\mid mh,\\
  2,&p\nmid mh.
 \end{cases}
\]
If $m$ and $h$ are both odd, then $n$ and $mn+h$ cannot both be odd primes
in the ranges $n\asymp_W X/m$ and $mn+h\asymp_W X$ once $X$ is large; the
prime--prime contribution is then zero.  In the admissible cases the upper
singular factor satisfies
\begin{equation}\label{eq:upper-singular-factor}
 \cS^+(m,h)
 \ll_h\prod_{\substack{p\mid m\\p>2}}\frac{p-1}{p-2}
 \ll_h\frac{m}{\varphi(m)}.
\end{equation}
Indeed, after comparison with $p/(p-1)$ the quotient of the two local
factors is $1+O(p^{-2})$.  The standard dimension-two upper sieve
\citep[Chapter~6]{IwaniecKowalski2004} therefore gives
\begin{equation}\label{eq:two-prime-row-upper}
 \sum_{\substack{n\geq1\\n,\,mn+h\ \mathrm{prime}}}
 (\log n)\log(mn+h)|W(mn/X)|
 \ll_{h,W}\frac{X}{\varphi(m)}.
\end{equation}
Summing \eqref{eq:two-prime-row-upper} and using
\eqref{eq:phi-harmonic-estimates} costs $O_{h,W}(X\log(2Y))$.

It remains to restore proper prime powers.  If $n$ is a proper prime
power, their aggregate contribution is bounded by
\[
 \sum_{m\leq Y}(X/m)^{1/2}(\log X)^2
 \ll \sqrt{XY}(\log X)^2.
\]
If $mn+h$ is a proper prime power, the crude aggregate bound
\[
 \sum_{m\leq Y}X^{1/2}(\log X)^2
 \ll YX^{1/2}(\log X)^2
\]
suffices.  Both are $o(X)$ in the stated range.  Finally, if $(m,h)>1$
and $\Lambda(mn+h)\ne0$, the same argument as in
\cref{lem:one-prime-edge} forces $mn+h=p^a$ for a fixed prime $p\mid h$;
the total is $X^{o(1)}$.  This proves the lemma.
\end{proof}

\begin{remark}[Why the power range is deliberately conservative]
The exponent $1/3$ is not optimized.  It makes all proper-prime-power
ledgers $o(X)$ with no need for a refined progression count.  The result
below concerns subpower hard edges, for which this range is more than
sufficient.
\end{remark}

\subsection{The bulk theorem}

\begin{theorem}[Two-sided hard-edge bulk persistence]
\label{thm:bulk-persistence}
Fix $h\ne0$ and $W\in C_c^1((0,\infty))$.  Uniformly for
\begin{equation}\label{eq:bulk-Y-range}
 2\leq Y\leq X^{1/3},
\end{equation}
one has
\begin{align}
 \cP_h^{\mathrm{bulk}(Y)}(X;W)
 &=I_0(W)X\log X+O_{h,W}(X\log(2Y)),
 \label{eq:prime-source-bulk}\\
 \cA_h^{\mathrm{bulk}(Y)}(X;W)
 &=(1-C_h)I_0(W)X\log X+O_{h,W}(X\log(2Y)).
 \label{eq:residual-bulk}
\end{align}
\end{theorem}

\begin{proof}
Choose $0<\alpha<\beta$ with $\supp W\subset[\alpha,\beta]$.  In the
range \eqref{eq:bulk-Y-range}, the intersection $m,n\leq Y$ is disjoint
from $mn\asymp_W X$ once $X$ is sufficiently large.  Thus each bulk form
is its complete form minus the $n\leq Y$ edge minus the $m\leq Y$ edge.

For the prime-source form, the $n$-edge is bounded by the
$\Lambda(q)$ part of \eqref{eq:one-prime-q-edge}; the $m$-edge is
\cref{lem:two-linear-forms-edge}.  Insert
\eqref{eq:prime-source-complete-asymptotic}.

For the residual form, bound $|\Lambda(n)-1|\leq\Lambda(n)+1$.  The
$n$-edge follows from \eqref{eq:one-prime-q-edge}.  On the $m$-edge, the
$\Lambda(n)$ part is \cref{lem:two-linear-forms-edge}, and the constant
part is \eqref{eq:one-prime-k-edge}.  Insert
\eqref{eq:residual-leading-only}.  All discarded $O(X)$ terms are absorbed
by $O(X\log(2Y))$.
\end{proof}

\begin{corollary}[Subpower edge removal]
\label{cor:subpower-edge-removal}
Suppose $2\leq Y=Y(X)\leq X^{1/3}$ and
$\log(2Y)=o(\log X)$.  If $I_0(W)\ne0$, then
\[
 \cP_h^{\mathrm{bulk}(Y)}(X;W)\sim I_0(W)X\log X.
\]
If additionally $(1-C_h)I_0(W)\ne0$, then
\[
 \cA_h^{\mathrm{bulk}(Y)}(X;W)
 \sim(1-C_h)I_0(W)X\log X.
\]
In particular, the conclusion applies to
$Y=(\log X)^K$ and to $Y=\exp((\log X)^\alpha)$ for fixed $K>0$ and
$0<\alpha<1$.
\end{corollary}

\begin{remark}[What persists and what does not]
When $Y=X^\delta$, the error in
\cref{thm:bulk-persistence} has the same order as the displayed main term;
no unchanged coefficient is then asserted.  Even in the subpower range,
the theorem controls a wide interior aggregate, not a fixed-width window
in $\log(n/m)$.  It is therefore neither a Type~II estimate nor a
prime-pair asymptotic.
\end{remark}
